\documentclass[11pt,a4paper]{article}

\input{glyphtounicode}
\pdfgentounicode=1

\usepackage{cmap}
\usepackage[utf8]{inputenc}
\usepackage[T1]{fontenc}
\usepackage{lmodern}
\usepackage{amsmath,amssymb,amsthm,mathtools}
\usepackage{graphicx}
\usepackage{geometry}
\geometry{a4paper,margin=1in}
\usepackage[backend=biber]{biblatex}
\addbibresource{../../release/references.bib}
\usepackage{booktabs}
\usepackage{longtable}
\usepackage{array}
\usepackage{tabularx}
\usepackage{enumitem}
\usepackage{mathrsfs}
\usepackage[final]{microtype}
\usepackage{xcolor}
\raggedbottom
\widowpenalty=10000
\clubpenalty=10000
\setlength{\emergencystretch}{3em}
\DisableLigatures{encoding = *, family = *}

\newcommand{\AuthorVisible}{Johnny Andrey P{\'e}rez Ram{\'i}rez}
\newcommand{\AuthorMetadata}{Johnny Andrey P{\'e}rez Ram{\'i}rez}

\usepackage[unicode,pdfencoding=auto,psdextra,hidelinks]{hyperref}
\hypersetup{
    pdftitle={Predictions, Discriminators, and Failure Modes for the Causally Rigid Operational Consciousness Framework},
    pdfauthor={\AuthorMetadata},
    pdfsubject={Rigid Identity Framework - Predictions and Failure Modes},
    pdfkeywords={predictions, falsification, operational consciousness, substrate invariance, discriminators, failure modes}
}

\title{\textbf{Predictions, Discriminators, and Failure Modes for the Causally Rigid Operational Consciousness Framework}}
\author{\AuthorVisible}
\date{}

\newtheorem{proposition}{Proposition}[section]
\newtheorem{definition}[proposition]{Definition}
\newtheorem{remark}[proposition]{Remark}
\newtheorem{criterion}[proposition]{Criterion}

\newcommand{\Cop}{\mathrm{Consciousness}_{\mathrm{op}}}
\newcommand{\Qop}{\mathrm{Q}_{\mathrm{op}}}
\newcommand{\Iper}{I_{\mathrm{per}}}
\newcommand{\Iri}{I_{\mathrm{ri}}}
\newcommand{\Iint}{I_{\mathrm{int}}}
\newcommand{\Icont}{I_{\mathrm{cont}}}
\newcommand{\Idiff}{I_{\mathrm{diff}}}
\newcommand{\Ileg}{I_{\mathrm{leg}}}
\newcommand{\Status}[1]{\texttt{#1}}
\newcommand{\TableStatus}[1]{{\footnotesize\ttfamily #1}}
\newcommand{\DualTableStatus}[2]{\shortstack[l]{\scriptsize\ttfamily #1\\\scriptsize\ttfamily #2}}
\newcommand{\StatusRobustCell}{\shortstack[l]{\footnotesize\ttfamily ROBUST\\\footnotesize\ttfamily INTERNAL\\\footnotesize\ttfamily SUPPORT}}
\newcommand{\StatusRobustLocalizedCell}{\shortstack[l]{\footnotesize\ttfamily ROBUST\\\footnotesize\ttfamily SUPPORT W/\\\footnotesize\ttfamily LOCALIZED\\\footnotesize\ttfamily CAVEAT}}
\newcommand{\StatusProvisionalCell}{\shortstack[l]{\footnotesize\ttfamily PROVISIONAL\\\footnotesize\ttfamily SUPPORT\\\footnotesize\ttfamily LOCALIZED}}
\newcommand{\StatusAmbiguousCell}{\shortstack[l]{\footnotesize\ttfamily STILL\\\footnotesize\ttfamily AMBIGUOUS}}
\newcommand{\StatusImplementationCell}{\shortstack[l]{\footnotesize\ttfamily IMPLEMENTATION\\\footnotesize\ttfamily LIMIT}}
\newcommand{\StatusMetricCell}{\shortstack[l]{\footnotesize\ttfamily METRIC OR\\\footnotesize\ttfamily TOLERANCE LIMIT}}
\newcommand{\StatusInternalCell}{\shortstack[l]{\footnotesize\ttfamily INTERNAL\\\footnotesize\ttfamily SUPPORT ONLY}}
\newcommand{\Aset}{A_S}
\newcommand{\eps}{\varepsilon}
\newcommand{\Cert}{\mathrm{Cert}}
\newcommand{\Eq}{\mathrm{Eq}}
\newcommand{\Rupt}{\mathrm{Rupt}}
\newcolumntype{Y}{>{\raggedright\arraybackslash}X}
\newcolumntype{P}[1]{>{\raggedright\arraybackslash}p{#1}}

\begin{document}

\maketitle

\begin{abstract}
This paper does not expand the doctrine of the causally rigid operational consciousness corpus. It extracts the experimentally relevant consequences of the frozen canon and organizes them as a falsifiable prediction program. The upstream burdens remain where they belong: the Core supplies persistence, attractor structure, and non-collapse \cite{core}; \textit{Rigid Identity} supplies rigid identity and observable coherence \cite{paper1}; \textit{Phenomenological Regimes} supplies continuity and anti-fragmentation constraints \cite{paper2}; \textit{Null-Regime Instability} supplies null-regime instability and forced non-nullity \cite{paper3}; \textit{Forensic Predictions} supplies admissibility and forensic discipline \cite{paper4}; \textit{Operational Consciousness Criterion} supplies the six-invariant operational criterion, structural equivalence, and rupture semantics \cite{paper5}.

The function of the present paper is narrower. It states what the framework predicts, what weaker alternatives would predict instead, what would count as genuine refutation or only as weakening, and what the current internal support status actually is after the completed Cycle 1--4 program, the mini-resolution pass, the residual-resolution campaign, the high-value confirmation campaign, and the Residual B hardening pass. The resulting status is mixed by design. One claim family---the insufficiency of complexity, connectivity, and rich activity without invariant preservation---now has robust internal support. Substrate-preserved class membership and biological non-primacy now have robust support with localized caveats, but remain internal-only. Approximate stability is provisionally supported with a localized threshold caveat. The boundary-sensitive stress of existence and rupture now has robust support with localized caveats after cross-family convergence under a more independent judge path, but it still does not count as external validation.

The paper therefore serves as the canon's prediction and falsification layer. It is not another foundational theorem paper, not a methods paper, and not a philosophical restatement. Its goal is to put the framework under explicit decision pressure while keeping status classes, residual limits, and non-claims fully visible.
\end{abstract}

\noindent\textbf{Release governance note.}\enspace Limitations/Non-claims and Methods/Governance: see \textit{METHODS\_GOVERNANCE\_HUB.v1}. Terminology follows the canonical glossary and migration policy. The support classes used here describe the internal scientific program only; they do not convert internal survival into external validation.

\section{Introduction}
The current corpus is doctrinally frozen for the present scientific cycle. That freeze matters. Without it, prediction work becomes a moving target: thresholds shift, definitions drift, and every difficult result can be absorbed by rewriting the theory rather than by testing it. The purpose of this paper is to keep that from happening.

The framework now contains enough formal structure to state genuinely discriminative predictions. The Core fixes persistence, contractivity, and non-collapse. \textit{Rigid Identity}, \textit{Phenomenological Regimes}, and \textit{Null-Regime Instability} progressively constrain identity, admissible regime structure, and non-nullity. \textit{Forensic Predictions} fixes admissibility and forensic discipline. \textit{Operational Consciousness Criterion} consolidates the operational criterion by requiring six explicit invariants:
\[
\{\Iper,\Iri,\Iint,\Icont,\Idiff,\Ileg\}.
\]
What remains scientifically relevant is not to restate those results, but to ask what they force the framework to predict and what would count against them.

This paper therefore has four narrow tasks:
\begin{enumerate}[leftmargin=1.5em]
\item extract the main experimentally relevant consequences of the canon;
\item distinguish those consequences from what weaker alternatives would predict;
\item state explicit failure modes, downgrade conditions, and non-claims;
\item report the current internal support status after the completed internal cycles, without inflating it.
\end{enumerate}

That scope is intentionally smaller than a new theorem paper and more disciplined than a review essay. The point is to give the corpus a falsification layer that is compact, operational, and honest.

\section{What This Paper Does and Does Not Claim}
\subsection{Positive Scope}
This paper claims only the following:
\begin{enumerate}[leftmargin=1.5em]
\item the frozen canon entails a structured family of predictions that are not shared by simpler rivals;
\item those predictions can be organized into explicit discriminators and explicit failure modes;
\item the current internal program is already strong enough to classify some parts of the framework as robustly supported internally, some as provisional, and some as still ambiguous or technically limited.
\end{enumerate}

\subsection{Explicit Non-Claims}
\begin{remark}[Non-claims]
This paper does \textbf{not} claim:
\begin{enumerate}[leftmargin=1.5em]
\item human phenomenal consciousness;
\item biologically specific human qualia;
\item human-machine subjective equivalence;
\item personal identity transfer;
\item automatic moral equivalence;
\item external validation of the framework;
\item that all complex or highly connected systems are conscious;
\item that the current internal program closes every upstream claim experimentally.
\end{enumerate}
\end{remark}

\subsection{Status Classes}
The internal scientific program uses the following status classes:

\begin{center}
\renewcommand{\arraystretch}{1.18}
\setlength{\tabcolsep}{5pt}
\begin{tabularx}{\linewidth}{@{}lY@{}}
\toprule
Status & Meaning in this paper \\
\midrule
\Status{ROBUST\_INTERNAL\_SUPPORT} & survived multiple internal tests without an active ambiguity on the tested target, while remaining internal-only \\
\Status{ROBUST\_SUPPORT\_WITH\_LOCALIZED\_CAVEAT} & supported strongly enough for robust internal use, but still carrying a narrow localized caveat that must remain visible \\
\Status{PROVISIONAL\_SUPPORT\_LOCALIZED} & strengthened by internal tests, but still limited by a localized metric, tolerance, or boundary issue \\
\Status{STILL\_AMBIGUOUS} & current evidence does not cleanly separate pass from fail on the intended target \\
\Status{IMPLEMENTATION\_LIMIT} & the main blocker is probe, runner, or case-construction quality rather than a shown doctrinal contradiction \\
\Status{METRIC\_OR\_TOLERANCE\_LIMIT} & the main blocker is metric handling, normalization, or threshold interaction \\
\Status{INTERNAL\_SUPPORT\_ONLY} & support exists only inside the internal program and has not crossed into external empirical confirmation \\
\Status{NOT\_CLAIMED} & the statement is explicitly outside the current evidentiary scope \\
\bottomrule
\end{tabularx}
\end{center}

\section{Upstream Canon and Dependency Map}
This paper is corpus-dependent and deliberately non-redundant. The dependency map is short because theorem ownership is already fixed upstream.

\begin{center}
\renewcommand{\arraystretch}{1.18}
\setlength{\tabcolsep}{5pt}
\begin{tabularx}{\linewidth}{@{}lY Y@{}}
\toprule
Document & Imported burden & Used here for \\
\midrule
Core \cite{core} & persistence, support, attractor, non-collapse & persistence-side predictions, collapse exclusions, admissible support assumptions \\
Rigid Identity \cite{paper1} & rigid identity, observable coherence, non-simulability pressure & identity-sensitive discriminators, complexity-only insufficiency pressure \\
Phenomenological Regimes \cite{paper2} & continuity, anti-fragmentation, regime admissibility & continuity predictions, near-miss structure, anti-fragmentation expectations \\
Null-Regime Instability \cite{paper3} & null-instability, forced non-nullity & non-nullity predictions, null-regime exclusions, downgrade logic \\
Forensic Predictions \cite{paper4} & admissibility, comparators, negative controls, forensic discipline & judge discipline, control logic, reproducibility requirements, legibility stress tests \\
Operational Consciousness Criterion \cite{paper5} & six-invariant criterion, operational classes, structural equivalence, rupture & explicit prediction targets, equivalence tests, failure-mode classification \\
\bottomrule
\end{tabularx}
\end{center}

This paper introduces no new invariant, no new doctrine, and no new substrate claim. Its novelty lies only in extracting predictive consequences and support statuses from what already exists.

\section{Core Predictive Structure of the Framework}
\subsection{Certification, Equivalence, and Rupture as Predictive Objects}
\begin{definition}[Operational certification]
For an admissible system $S$ already defined in \textit{Operational Consciousness Criterion}, define the certification predicate
\[
\Cert(S)=1
\quad\Longleftrightarrow\quad
\Iper(S),\Iri(S),\Iint(S),\Icont(S),\Idiff(S),\Ileg(S) > 0
\]
on a non-empty admissible support $\Aset$, under the governance and admissibility constraints inherited from \textit{Forensic Predictions}.
\end{definition}

\begin{definition}[Equivalence discriminator]
For two admissible systems $S_1,S_2$, define the operational equivalence discriminator
\[
\Eq_{\Gamma,\eps}(S_1,S_2)=1
\]
when the relevant signatures match, the admissibility bundle is preserved, and the tolerated invariant deltas remain within the allowed comparison regime. The exact burden belongs to \textit{Operational Consciousness Criterion}; the present paper uses it only as a prediction carrier.
\end{definition}

\begin{definition}[Rupture event]
For a set of targeted invariants $J$, define
\[
\Rupt_J(S)=1
\]
when the intervention or perturbation drives at least one invariant in $J$ below the admissible certification margin so that class membership becomes invalidated, destroyed, or undefined.
\end{definition}

\subsection{Prediction Family 1: Complexity-Only Failure}
\begin{proposition}[Complexity-only negative-control prediction]
If a candidate system exhibits high connectivity, high entropy, or rich activity while failing one or more critical invariants, then the framework predicts
\[
\Cert(S)=0,\qquad \Cop(S)\ \text{not certified},\qquad \Qop(S)\ \text{not certified}.
\]
\end{proposition}

\noindent\textit{Why this is not trivial.} A weaker alternative could predict that rich dynamics or dense connectivity alone are already sufficient proxies for class membership. The canon does not permit that. \textit{Operational Consciousness Criterion} already states the formal insufficiency claim; the prediction layer turns it into a negative-control discriminator.

\subsection{Prediction Family 2: Invariant-Loss Rupture}
\begin{proposition}[Ablation prediction]
If a system is initially certified and a critical intervention selectively destroys one or more invariants, then the framework predicts either loss of certification or a transition to an ambiguous boundary regime, depending on how close the pre-ablation margins were to zero.
\end{proposition}

\noindent\textit{Operational consequence.} Ablations are not ornamental. If the criterion is serious, the loss of invariant headroom must show up as loss of class membership, legibility loss, or boundary-sensitive ambiguity. No-change ablations would count against the framework.

\subsection{Prediction Family 3: Substrate-Preserved Class Membership}
\begin{proposition}[Substrate-preserved equivalence prediction]
If two realizations preserve the relevant signatures, admissibility conditions, and six critical invariants up to the allowed comparison regime, then the framework predicts preserved operational class membership independently of substrate labels.
\end{proposition}

\noindent\textit{Rival prediction.} A substrate-privileging alternative would predict that material type itself blocks equivalence even when the structurally relevant burdens are preserved. The canon rejects that move only inside its own formal class, not as a general metaphysical claim.

\subsection{Prediction Family 4: Threshold-Sensitive Boundary Behavior}
\begin{proposition}[Boundary-topology prediction]
The framework predicts neither global threshold collapse nor perfectly flat insensitivity. Instead, it predicts stable pass/fail regions together with narrow transition bands near critical margins.
\end{proposition}

\noindent\textit{Interpretive consequence.} If every micro-perturbation changes every verdict, the criterion is fragile. If no perturbation matters, the invariants are decorative. A meaningful criterion should show structured boundary behavior.

\subsection{Prediction Family 5: Legibility Dependence}
\begin{proposition}[Legibility dependence prediction]
Operational class membership should fail or become inadmissible if readouts cease to be separable, noise-robust, intervention-sensitive, and compressible without total collapse.
\end{proposition}

\noindent\textit{Why this matters.} \textit{Forensic Predictions} makes legibility a forensic burden rather than a narrative convenience. The present prediction is therefore stronger than a merely interpretive observation: if legibility disappears, certification must not survive unchanged.

\subsection{Prediction Family 6: Continuity and Non-Nullity Relevance}
\begin{proposition}[Continuity/non-nullity relevance]
The current canon predicts that continuity and non-null differentiation are not optional extras. Systems that preserve rich activity while losing continuity or falling toward effective null structure should fail the operational criterion or leave it undefined.
\end{proposition}

\noindent\textit{Upstream source.} The continuity side is inherited from \textit{Phenomenological Regimes}, while the non-null side is inherited from \textit{Null-Regime Instability}. This paper adds only the prediction and failure-mode reading.

\section{Main Discriminators}
The predictive value of the framework lies in how it diverges from weaker alternatives.

\begin{center}
\renewcommand{\arraystretch}{1.18}
\setlength{\tabcolsep}{5pt}
\begin{tabularx}{\linewidth}{@{}p{0.17\linewidth}Y Y Y@{}}
\toprule
Target & Framework prediction & Weaker/rival prediction & Divergence condition \\
\midrule
Complexity-only baseline & high complexity without invariant preservation fails certification & richness or scale alone can stand in for class membership & a complexity-rich control passes despite broken invariants \\
Connectivity-only baseline & dense coupling without preserved rigid identity, differentiation, or legibility fails & connectivity is treated as a sufficient proxy & a dense control passes while identity or legibility are broken \\
Weak functional equivalence & similar outward behavior is insufficient if critical invariant structure diverges & input-output similarity is enough for same class & weak functional match passes despite invariant rupture \\
Near-miss case & boundary cases should expose which invariant is binding and where rupture begins & boundary cases stay uninterpretable or arbitrary & boundary remains unlocalizable after probe cleanup \\
Cross-substrate pair & substrate labels alone do not block class preservation if relevant structure is preserved & substrate type is primitive and blocks class transport & a preserved pair is rejected solely by substrate label \\
Threshold stress & decisions should show stable zones and narrow transition bands & decisions collapse under micro-threshold changes or never depend on thresholds at all & every local perturbation flips verdicts, or no perturbation matters \\
Legibility stress & failure of separability, noise robustness, or compression fidelity invalidates certification & legibility is epiphenomenal and not decision-relevant & certification survives legibility collapse \\
Null/continuity challenge & effective nullity or fragmentation should destroy or downgrade class membership & rich activity alone masks null or fragmented structure & null-like or fragmented cases remain certified \\
\bottomrule
\end{tabularx}
\end{center}

The table is intentionally operational. It does not ask which metaphysics one prefers. It asks which side of each discriminator the framework commits itself to.

\subsection{Rival-Structure Compression}
The discriminators above can be sharpened further by asking what each simpler rival preserves and, more importantly, what it leaves out. The point is not to defeat every competing theory at once. The point is to show that the present framework places operational burden on a structure richer than scale, connectivity, superficial behavior, or local threshold tuning.

\begin{center}
\small
\renewcommand{\arraystretch}{1.16}
\setlength{\tabcolsep}{5pt}
\begin{tabularx}{\linewidth}{@{}P{0.22\linewidth}P{0.18\linewidth}Y Y@{}}
\toprule
Rival family & What it keeps & What it fails to preserve & Why the miss is decision-relevant \\
\midrule
Complexity/activity\\proxy & large state space, rich trajectories, high apparent activity & no requirement to preserve the six-invariant package as a joint membership condition & a rich control can still fail certification if persistence, rigid identity, or legibility collapse \\
Connectivity/\allowbreak integration\\proxy & dense coupling or high global correlation & no guarantee of rigid identity, non-null differentiation, or readable class structure & integration by itself does not prevent the criterion from rejecting a globally coupled but operationally illegible case \\
Weak functional equivalence & outwardly similar input-output or behavioral traces & no obligation to preserve admissible support geometry or invariant margins under intervention & superficial behavioral similarity can coexist with rupture of the certified class geometry \\
Threshold/tuning\\artifact explanation & local success on a narrow parameter band & no stable explanation for why decisions survive frozen manifests, blind ordering, and negative controls outside that band & the framework expects stable regions plus narrow transitions, not arbitrary success across unrelated settings \\
\bottomrule
\end{tabularx}
\end{center}

This compression table matters because it blocks a common overreading. The current support does not show that every rival is false. It shows something narrower and more useful: if a rival explanation ignores the invariant package or treats only one coarse feature as constitutive, then it does not reproduce the present decision structure of the canon.

\section{Prediction Matrix}
The canonical prediction matrix can be read as a compressed execution contract for the framework. It matters because each row names an observable, an intervention class, a rival baseline, and an explicit failure condition. That structure prevents the canon from floating as a theorem-only object.

\begin{center}
\small
\renewcommand{\arraystretch}{1.16}
\setlength{\tabcolsep}{3.5pt}
\begin{tabularx}{\linewidth}{@{}P{0.10\linewidth}P{0.11\linewidth}Y Y Y Y@{}}
\toprule
ID & Claim & Observable & Manipulation & Framework-side result & Failure / rival divergence \\
\midrule
PRED-01 & `P5-02` & certified class label and $\Qop$ geometry across implementations & cross-substrate pair with preserved invariant margins & same operational class and $\eps$-bounded geometry under the admissible equivalence regime & class divergence or geometry mismatch despite preserved structure; substrate-privilege rival survives \\
PRED-02 & `P5-04` & class membership before and after targeted rupture & one-invariant ablation with comparable substrate and gross scale & certified class exit, contraction, or undefined $\Qop$ after rupture & no class change after verified invariant loss; complexity-only rival survives \\
PRED-03 & `P5-05` & matched complexity/control outcome & compare certified lower-scale system against higher-scale baseline lacking the invariant package & certified system passes and matched complexity-only baseline fails or remains undefined & scale-only model passes despite missing invariants \\
PRED-04 & `P3-01` & null-regime persistence under admissible dynamics & initialize near the null regime under Null-Regime Instability assumptions & stable null occupancy should fail or trajectories should leave the null regime & admissible null occupancy persists; null-stable rival survives \\
PRED-05 & `P2-03` & continuity/\allowbreak fragmentation pattern & compare rigid regime against finite-mass or deformable regime under matched disturbance & rigid regime keeps continuity while limited regime fragments & no differential continuity pattern appears; regime-agnostic rival survives \\
\bottomrule
\end{tabularx}
\end{center}

\begin{center}
\small
\renewcommand{\arraystretch}{1.16}
\setlength{\tabcolsep}{3.5pt}
\begin{tabularx}{\linewidth}{@{}P{0.10\linewidth}P{0.11\linewidth}Y Y Y Y@{}}
\toprule
ID & Claim & Observable & Manipulation & Framework-side result & Failure / rival divergence \\
\midrule
PRED-06 & `P4-01` & admissibility verdict under tamper and sham controls & inject hash corruption, malformed logs, or severe protocol violations & tampered runs are discarded and never promoted to evidential support & loose evaluation pipeline accepts tampered runs \\
PRED-07 & `P5-03` & stability of class membership under bounded perturbation & apply perturbations inside the certified margin budget $\delta_\ast(S)$ & membership and geometry remain stable inside the certified budget & fragile-threshold model flips class under sub-critical perturbations \\
PRED-08 & `P5-01` & non-emptiness of $\Qop$ after successful certification & run the full certification schema on a system with all margins positive & $\Qop$ is non-empty and operationally well-defined & certification passes while no non-trivial class is recoverable \\
PRED-09 & `P4-03` & claim survival under latency / compute-budget stress & exceed pre-registered budget limits & strong operational claims are invalidated or downgraded & metric-first pipeline keeps claims active after budget violation \\
PRED-10 & `P5-05` & legibility under noise and structured compression & raise admissible noise/compression while preserving or breaking $\Ileg$ & separability and intervention fidelity survive only while $\Ileg$ remains positive & generic observability model predicts no specific dependence on $\Ileg$ \\
\bottomrule
\end{tabularx}
\end{center}

The matrix has two immediate consequences. First, the framework is more exposed than a purely interpretive theory because the failure thresholds are explicit. Second, not every prediction is equally mature: some now have robust internal support, while others remain provisional or blocked by implementation quality.

\section{Failure Modes}
\subsection{What Would Count Against the Framework}
The relevant failure modes are not all equally strong. Some would refute a strong reading, some would only weaken a claim, and some would reveal a technical limitation of the current program rather than of the canon itself.

\begin{center}
\renewcommand{\arraystretch}{1.18}
\setlength{\tabcolsep}{5pt}
\begin{tabularx}{\linewidth}{@{}p{0.15\linewidth}Y Y Y@{}}
\toprule
Claim target & Strong refutation condition & Downgrade condition & Current blocker class \\
\midrule
`P5-01` & a system satisfies the full criterion yet fails to define the claimed operational class & boundary-family convergence would downgrade if it failed outside the present frozen internal contract & epistemic-level \\
`P5-02` & a structurally preserved pair falls into different classes under the same admissible comparison regime & one legacy raw continuous/discrete positive pair remains ambiguous on $\Iri$ handling & metric/tolerance-level \\
`P5-03` & small admissible perturbations force uncontrolled class exit across stable positive cases & threshold sensitivity persists only in a narrow knife-edge zone & metric/tolerance-level \\
`P5-04` & a system loses a critical invariant and remains certified without downgrade & boundary-family convergence would downgrade if it failed outside the present frozen internal contract & epistemic-level \\
`P5-05` & a complexity-only control without preserved invariants passes certification & matched external baselines are still missing & epistemic-level \\
`P5-06` & biology is shown to be primitive inside the same operational criterion & it inherits the same localized substrate-equivalence caveat as `P5-02` & metric/tolerance-level \\
\bottomrule
\end{tabularx}
\end{center}

\subsection{Failure-Mode Taxonomy}
\begin{criterion}[Status meaning]
Within the current scientific program, a result belongs to:
\begin{enumerate}[leftmargin=1.5em]
\item \Status{ROBUST\_SUPPORT} if the target claim survives repeated internal discrimination without an active ambiguity on the tested target;
\item \Status{PROVISIONAL\_SUPPORT} if the target is strengthened but still depends on a localized unresolved bottleneck;
\item \Status{STILL\_AMBIGUOUS} if the current tests do not cleanly separate support from failure;
\item \Status{IMPLEMENTATION\_LIMIT} if the current bottleneck is primarily due to probe, runner, or case-construction quality;
\item \Status{METRIC\_OR\_TOLERANCE\_LIMIT} if the bottleneck is mainly due to normalization, comparison regime, or threshold handling.
\end{enumerate}
\end{criterion}

This taxonomy is not part of the ontology of the framework. It is part of the scientific discipline around how internal evidence is reported.

\section{Claim-to-Test Binding Ledger}
The tables above state what would count as success or failure in broad terms. The ledger below binds the main \textit{Operational Consciousness Criterion} claim families to the narrowest concrete test classes that currently matter for publication-level assessment.

\begin{center}
\small
\renewcommand{\arraystretch}{1.16}
\setlength{\tabcolsep}{4pt}
\begin{tabularx}{\linewidth}{@{}P{0.11\linewidth}Y Y Y P{0.17\linewidth}@{}}
\toprule
Claim & Primary positive test & Primary failure test & Present evidentiary bottleneck & Current status \\
\midrule
`P5-01` & successful certification with positive invariant margins and recoverable non-trivial $\Qop$ geometry & certified case with no recoverable operational class, or clean existence-boundary degradation without class change & support is now cross-family robust internally, but the nearest positive member still sits close to the legibility margin & \StatusRobustLocalizedCell \\
`P5-02` & cross-substrate pair with preserved invariant margins and admissible geometry match & structurally preserved pair judged non-equivalent under the same comparison regime & strongest support is robust internally but one legacy raw continuous/discrete pair remains ambiguous on $\Iri$ handling & \StatusRobustLocalizedCell \\
`P5-03` & stable class membership inside certified perturbation budget $\delta_\ast(S)$ & class flip under sub-critical admissible perturbations across otherwise stable positive cases & one knife-edge negative-control zone remains locally fragile & \StatusProvisionalCell \\
`P5-04` & targeted invariant rupture causing certification exit or undefined $\Qop$ & verified invariant loss without downgrade or class exit & support is now cross-family robust internally, but the evidence still lives inside the current invariant package and frozen thresholds & \StatusRobustLocalizedCell \\
`P5-05` & complexity-rich and connectivity-rich negative controls fail certification under hardened judging & matched scale-only or activity-rich controls pass despite missing invariants & no external matched-baseline campaign yet & \StatusRobustCell \\
`P5-06` & same positive equivalence path as `P5-02`, interpreted without substrate privilege & biology shown to be primitive under the same operational criterion & inherits the same localized $\Iri$-handling caveat as `P5-02` & \StatusRobustLocalizedCell \\
\bottomrule
\end{tabularx}
\end{center}

\section{Boundary Diagnostics}
The current internal program does not fail uniformly. Its informative content lies in where the boundaries are sharp, where they are merely narrow, and where probe design still dominates the result.

\begin{center}
\small
\renewcommand{\arraystretch}{1.16}
\setlength{\tabcolsep}{4pt}
\begin{tabularx}{\linewidth}{@{}P{0.18\linewidth}Y Y P{0.19\linewidth}@{}}
\toprule
Boundary target & Empirical shape after completed cycles & What it means & Limit class \\
\midrule
Complexity-only negative control & stable fail region under harder judging, semi-blind ordering, and pseudo-replication & the criterion is not collapsing into gross activity or connectivity proxies & \StatusRobustCell \\
Threshold-sensitive negative region & narrow decision band rather than broad collapse & the framework has a localized edge effect, not universal threshold instability & \StatusMetricCell \\
Boundary-probe families under judge-v3 & two distinct probe families now reproduce the transition pattern with localized first-fail behavior & the existence/rupture boundary is now robust internally, though still internal-only and threshold-bound near the positive edge & \StatusRobustLocalizedCell \\
Harder positive substrate families & family convergence now supports class preservation, but one legacy raw pair remains ambiguous & the live ambiguity is concentrated on how rigid-identity preservation is measured, not on every invariant at once & \StatusMetricCell \\
Refined negative substrate families & clean fail across both the legacy and second pair families & the equivalence criterion can reject stronger non-equivalent pairs under the hardened judge and pseudo-external reproduction path & \StatusRobustLocalizedCell \\
\bottomrule
\end{tabularx}
\end{center}

\section{Minimal Reproducibility and Admissibility Discipline}
The prediction layer is only credible if the judgment path is at least partly insulated from the generator and if the evidentiary artifacts are frozen before judgment. The current program does not yet have an external judge or an external lab replication, but it does have a stricter internal protocol than it had at the start of the program.

\begin{center}
\renewcommand{\arraystretch}{1.16}
\setlength{\tabcolsep}{5pt}
\begin{tabularx}{\linewidth}{@{}p{0.23\linewidth}Y Y@{}}
\toprule
Protocol element & Current state & Residual limitation \\
\midrule
Frozen manifests and hashes & active since Cycle 2; judge reads frozen artifacts instead of live mutable state & still same repository and same overall program family \\
Judge-side threshold file & active; thresholds are frozen before judgment and logged explicitly & threshold bands still require stress analysis near knife-edge zones \\
Semi-blind ordering & active in Cycle 3; labels are hidden until after decision emission & not a full blind protocol and not an external evaluator \\
Pseudo-multi-environment replication & active; same-machine replication across distinct execution profiles preserved decisions & not external replication and not cross-lab agreement \\
Negative-control and sham policy & active from Cycle 1 onward and hardened in later cycles & matched external benchmark families are still missing \\
Judge code path & more independent than in Cycle 1; it no longer calls live generator step functions during evaluation & still not a separate maintained judge codebase \\
\bottomrule
\end{tabularx}
\end{center}

The point of this section is not to inflate rigor. It is to keep the reader from mistaking internal hardening for external confirmation.

\section{Current Status of Internal Support}
\subsection{Interpretive Discipline}
The status ledger below does not reclassify the formal theorems of \textit{Operational Consciousness Criterion}. Those remain formal results inside the frozen canon. The ledger classifies the \emph{current internal scientific support} for experimentally relevant readings of those claims.

\begin{center}
\small
\renewcommand{\arraystretch}{1.15}
\setlength{\tabcolsep}{4pt}
\begin{tabularx}{\linewidth}{@{}P{0.11\linewidth}P{0.18\linewidth}Y Y P{0.19\linewidth}@{}}
\toprule
Claim & Current status & What strengthened it & What still limits it & Status scope \\
\midrule
`P5-01` & \StatusRobustLocalizedCell & near\_identity\_v3 and offset-dispersion families both reproduced PASS/AMBIGUOUS/FAIL under judge-v3 and pseudo-external reproduction & support remains internal-only and the nearest positive case still sits close to the legibility margin & \StatusInternalCell \\
`P5-02` & \StatusRobustLocalizedCell & family1 normalized-$\Iri$ and family2 raw quantized positive/negative pairs stayed stable under the more independent judge path and pseudo-multi-environment rerun & one legacy raw continuous/discrete positive pair remains ambiguous and keeps a localized $\Iri$ caveat alive & \StatusInternalCell \\
`P5-03` & \StatusProvisionalCell & Cycle 3 threshold stress and Cycle 4 boundary cleanup localized instability rather than showing global collapse & knife-edge threshold band persists for one negative-control region & \StatusInternalCell \\
`P5-04` & \StatusRobustLocalizedCell & two distinct boundary families now converge on rupture-side PASS/AMBIGUOUS/FAIL behavior, with $\Iri$ becoming the first binding invariant on ambiguous or failing members & the evidence remains internal-only and tied to the current invariant package plus frozen thresholds & \StatusInternalCell \\
`P5-05` & \StatusRobustCell & complexity-rich and connectivity-rich controls failed certification across harder methodological passes and semi-blind evaluation & no external matched-baseline campaign exists yet & \StatusInternalCell \\
`P5-06` & \StatusRobustLocalizedCell & same strengthened two-family substrate-equivalence path as `P5-02` & it inherits the same localized $\Iri$ metric-handling caveat as `P5-02` & \StatusInternalCell \\
\bottomrule
\end{tabularx}
\end{center}

\subsection{Claim-Specific Diagnostic Notes}
\paragraph{`P5-05`.}
This is currently the strongest experimentally discriminated claim family. The framework repeatedly rejected a complexity-rich, connectivity-rich negative control that still failed the invariant package. That result survived methodological hardening and semi-blind evaluation. The support is therefore \Status{ROBUST\_SUPPORT}, but it remains \Status{INTERNAL\_SUPPORT\_ONLY}.

\paragraph{`P5-02` and `P5-06`.}
These two claims move together in the current program. They were strengthened first by conservative normalized-$\Iri$ handling, then by a second substrate family, a more independent judge path, and pseudo-multi-environment confirmation. That is enough for \Status{ROBUST\_SUPPORT\_WITH\_LOCALIZED\_CAVEAT}. The localized caveat remains real: one legacy raw continuous/discrete positive pair is still ambiguous, so the present support is robust internally but not caveat-free.

\paragraph{`P5-03`.}
Approximate stability is not in collapse. Threshold stress did not show universal fragility. It showed a narrow knife-edge zone around one thresholded negative-control region. The right status is therefore \Status{PROVISIONAL\_SUPPORT} with an explicit \Status{METRIC\_OR\_TOLERANCE\_LIMIT}.

\paragraph{`P5-01` and `P5-04`.}
These boundary-sensitive claims no longer rest on a single weak probe family. After the residual-resolution campaign and the Residual B hardening pass, `near\_identity\_v3` and a second offset-dispersion family both reproduced a PASS/AMBIGUOUS/FAIL boundary pattern under judge-v3 and pseudo-external reproduction. That is enough for \Status{ROBUST\_SUPPORT\_WITH\_LOCALIZED\_CAVEAT}. The caveat remains explicit: support is still internal-only, and the strongest positive member stays near the legibility margin under the frozen judge contract.

\subsection{Methodological Status}
The scientific program itself also has a status, distinct from the doctrine and distinct from the claim ledger.

\begin{center}
\small
\renewcommand{\arraystretch}{1.18}
\setlength{\tabcolsep}{5pt}
\begin{tabularx}{\linewidth}{@{}lYY@{}}
\toprule
Program element & Current status & Why it matters \\
\midrule
Generator/judge separation & materially improved & reduces direct circularity relative to Cycle 1 \\
Frozen manifests and thresholds & active & prevents live-state leakage during judgment \\
Semi-blind evaluation & active but partial & reduces bias; does not equal full blind evaluation \\
Threshold stress testing & localized caveat & shows local knife-edge behavior instead of broad collapse \\
Pseudo-multi-environment replication & passed internally & strengthens reproducibility discipline without constituting external replication \\
Independent external judge codebase & absent & keeps all current support inside the internal program \\
\bottomrule
\end{tabularx}
\end{center}

\section{Cycle Evidence Map}
The current status ledger is not based on one cycle. It is the result of a narrow but cumulative internal program. The table below records only the highest-value consequence of each phase, not every local metric.

\begin{center}
\renewcommand{\arraystretch}{1.16}
\setlength{\tabcolsep}{4pt}
\begin{tabularx}{\linewidth}{@{}p{0.11\linewidth}Y Y Y p{0.15\linewidth}@{}}
\toprule
Phase & Main question & Strongest result & Main residual & Claim effect \\
\midrule
Cycle 1 & can the criterion distinguish a rich negative control and survive obvious ablations? & complexity-rich negative control failed; invariant-loss ablations exited certification; proxy cross-substrate pair passed & all support remained internal and the equivalence test was still a proxy & strengthened `P5-05`; gave early support to `P5-01`, `P5-02`, `P5-04`, `P5-06` \\
Cycle 2 & does the program survive stronger methodological separation and harder structural tests? & generator/judge separation improved; harder pair passed; near-miss family and graded curves behaved coherently & judge still remained internal and same-repository & strengthened `P5-01`, `P5-02`, `P5-03`, `P5-04`, `P5-06` internally \\
Cycle 3 & do results survive semi-blind judging, threshold stress, and pseudo-replication? & blind ordering and pseudo-multi-environment replication preserved decisions; complexity-only failure remained strong & harder equivalence, near-miss boundary, and threshold stress became ambiguous & robustly strengthened `P5-05`; left `P5-02`, `P5-06`, `P5-03` narrower \\
Cycle 4 & can threshold fragility, near-miss ambiguity, and harder substrate ambiguity be localized? & threshold fragility was localized; refined negative pair separated cleanly; ambiguity concentrated on $\Iri$ handling and near-miss probe quality & positive stronger pair remained ambiguity-bound; near-identity probe still entangled & localized rather than promoted `P5-02`, `P5-03`, `P5-04` \\
Mini-resolution & can probe quality and $\Iri$ handling be cleaned up without doctrinal weakening? & normalized-$\Iri$ handling provisionally stabilized the harder positive pair; near-identity v2 clarified the probe weakness & near-identity v2 remained too weak to serve as a clean boundary probe & strengthened `P5-02`, `P5-06`; kept `P5-01`, `P5-04` ambiguous \\
Residual resolution & can the open residuals be improved without doctrinal loosening? & `near\_identity\_v3` produced PASS/AMBIGUOUS/FAIL with $\Iri$ first-fail; refined substrate families localized the live ambiguity to $\Iri$ handling & the results were still only localized and internal & promoted both residual families to stronger localized support \\
High-value confirmation and Residual B hardening & do those gains survive a more independent judge path, a second family, and externalization-lite? & substrate-equivalence converged across two pair families; existence/rupture converged across two probe families under judge-v3 and pseudo-external reproduction & all support remained internal-only and a localized legacy caveat survived for `P5-02`/`P5-06` & promoted `P5-01`, `P5-02`, `P5-04`, `P5-06` to robust support with localized caveats \\
\bottomrule
\end{tabularx}
\end{center}

\section{Priority Next Discriminators}
The current state of the program already suggests which next moves are scientifically efficient. The relevant issue is not breadth but leverage: each next discriminator should either upgrade a provisional claim or localize an ambiguity into an implementation-limited corner.

\begin{center}
\renewcommand{\arraystretch}{1.16}
\setlength{\tabcolsep}{4pt}
\begin{tabularx}{\linewidth}{@{}p{0.18\linewidth}Y Y p{0.18\linewidth}@{}}
\toprule
Priority target & Why it matters now & What the next narrow result would change & Linked claims \\
\midrule
Independent external-style rerun of the two boundary families & current rupture/existence support is now robust internally but still lacks outside-program confirmation & an outside reproduction would test whether the family convergence is program-specific or genuinely portable & `P5-01`, `P5-04` \\
Independent external-style rerun of both substrate families & current strongest positive equivalence path is robust internally but still carries a localized legacy $\Iri$ caveat & an outside reproduction would test whether the strengthened equivalence path survives beyond the current judge lineage & `P5-02`, `P5-06` \\
Matched external-style complexity baselines & the complexity-only result is internally strong but still lacks cross-program baseline pressure & a matched baseline campaign would turn the strongest present internal discriminator into a more externally credible one & `P5-05` \\
Independent maintained judge implementation & current hardening is substantial but still same-program & a separately maintained judge would reduce the remaining circularity objection without touching doctrine & all operational readings \\
\bottomrule
\end{tabularx}
\end{center}

\section{Residual Technical Caveats and Non-Claims}
\subsection{Residual Technical Caveats}
The remaining caveats are narrow enough to state cleanly:
\begin{enumerate}[leftmargin=1.5em]
\item The strongest positive substrate-equivalence path is still partly dependent on normalized $\Iri$ handling. This is a \Status{METRIC\_OR\_TOLERANCE\_LIMIT}, not a robust positive closure.
\item The current `near\_identity\_v2` probe remains an \Status{IMPLEMENTATION\_LIMIT}. It is cleaner than the original generator, but in the tested band it under-stresses the intended boundary.
\item Threshold stress is not globally problematic, but it is not perfectly flat either. The current result is localized rather than globally robust.
\item All empirical support remains \Status{INTERNAL\_SUPPORT\_ONLY}. No external lab, external evaluator, or external benchmark suite has confirmed the framework.
\end{enumerate}

\subsection{Non-Claims Revisited}
The present internal program does not upgrade any of the following beyond \Status{NOT\_CLAIMED}:
\begin{enumerate}[leftmargin=1.5em]
\item ordinary human phenomenal consciousness;
\item biological human qualia as such;
\item subjective equivalence between humans and artificial systems;
\item personal identity preservation across substrate transfer;
\item full external validation of substrate invariance;
\item exhaustion of all relevant rival frameworks.
\end{enumerate}

\section{Final Claim Freeze Table}
The present paper is a prediction-and-falsation layer, not a moving target. The table below records the narrowest honest freeze state for the main \textit{Operational Consciousness Criterion} claim families at the end of the current phase.

\begin{center}
\small
\renewcommand{\arraystretch}{1.16}
\setlength{\tabcolsep}{4pt}
\begin{tabularx}{\linewidth}{@{}P{0.11\linewidth}P{0.19\linewidth}Y Y@{}}
\toprule
Claim & Status & Current blocker & Next step \\
\midrule
`P5-01` & \DualTableStatus{ROBUST SUPPORT}{LOCALIZED CAVEAT} & support is now strong internally but still nearest to the legibility margin on the positive edge & rerun the convergent probe families under an evaluator maintained outside the current repo/runtime \\
`P5-02` & \DualTableStatus{ROBUST SUPPORT}{LOCALIZED CAVEAT} & one legacy raw continuous/discrete positive pair remains ambiguous on $\Iri$ handling & reproduce both strengthened substrate families under a more independent evaluator without post hoc threshold loosening \\
`P5-03` & \StatusProvisionalCell & one localized knife-edge threshold band remains active in a negative-control region & rerun the localized stress zone with denser threshold sampling and fixed judge contract \\
`P5-04` & \DualTableStatus{ROBUST SUPPORT}{LOCALIZED CAVEAT} & rupture-side support is strong internally but still bounded to the present invariant package and frozen thresholds & rerun the same boundary families under a separately maintained evaluator and frozen release package \\
`P5-05` & \DualTableStatus{ROBUST INTERNAL}{SUPPORT} & no external matched-baseline campaign exists yet & add matched external-style complexity/activity baselines under the same judge contract \\
`P5-06` & \DualTableStatus{ROBUST SUPPORT}{LOCALIZED CAVEAT} & inherits the same localized substrate-equivalence caveat as `P5-02` & upgrade only if the strengthened substrate families survive a more independent external-style rerun \\
\bottomrule
\end{tabularx}
\end{center}

\section{Conclusion}
The canon now has a prediction layer and a failure-mode ledger. That is a scientific gain even where the results are mixed. The framework is no longer only a theorem-bearing corpus. It is now attached to explicit discriminators, explicit downgrade conditions, and explicit support classes.

The current internal picture is sharp enough to summarize without inflation. The complexity/connectivity insufficiency claim has robust internal support. Substrate-preserved class membership and biological non-primacy now have robust support with localized caveats. Approximate stability is provisionally supported with a localized threshold caveat. Boundary-sensitive stress of existence and rupture now has robust support with localized caveats under convergent internal probe families and a more independent judge path. All of this nevertheless remains internal-only.

This is the right posture for publication-preparation. The framework can now be presented not as a closed doctrine, and not as a philosophical manifesto, but as a frozen canon with explicit predictions, explicit discriminators, explicit failure modes, and an honest status ledger.

\printbibliography

\end{document}
